\subsection{Lifting the remaining projective-linear quotient through its full pronilpotent kernel}

\begin{statement}{D.42}{Proved}{An actual simultaneous relative-prime lift and its exact common quotient}
Retain the marked quotient $\mathcal W$ of D.41, its
automorphism $\alpha_{41}$, and its source kernel
\[
 R_W=R_T\cap K_{\mathcal R}\cap K_{\mathcal S}
       \lhd\widehat B .
\]
Let $\sigma:\widehat B\twoheadrightarrow
J=\operatorname{PGL}_2(\mathbb F_7)$ be the specific
map (D.41.6), and put $N_J=\ker(B\to J)$.
This is a free group of rank $57$. For every prime $p$,
let
\[
 R_{J,p}=\ker\bigl(\widehat N_J\longrightarrow
                              N_J^{(p)}\bigr),\qquad
 \mathcal P_{J,p}=\widehat B/R_{J,p},\qquad
 R_{J,\mathrm{nil}}=\bigcap_pR_{J,p}.                    \tag{D.42.1}
\]
Here $N_J^{(p)}$ is the full pro-$p$ completion of the
indicated marked free group. These are actual source
quotients, and
\[
 \mathcal P_{J,\mathrm{nil}}
   :=\widehat B/R_{J,\mathrm{nil}}
    =\mathop{\prod_J}_{p}\mathcal P_{J,p},\qquad
 \ker(\mathcal P_{J,\mathrm{nil}}\to J)
    =\prod_p N_J^{(p)}.                                \tag{D.42.2}
\]

There is an explicit maximal common marked quotient
of $\mathcal W$ and $\mathcal P_{J,\mathrm{nil}}$.
Use D.37's source map $B\to H$, $|H|=6912$, and set
\[
 H_{54}=H/O_2(H),\quad |O_2(H)|=128,\quad
 L_{54}=\ker(B\to H_{54}),\qquad
 \mathcal C_{54}
   =\varprojlim_{n\ge1}B/[L_{54},L_{54}]L_{54}^{3^n}.
                                                               \tag{D.42.3}
\]
The inverse system uses the actual quotient maps.
Its structure is
\[
 |H_{54}|=54,\qquad
 \operatorname{rank}L_{54}=10,\qquad
 1\longrightarrow\mathbb Z_3^{10}
  \longrightarrow\mathcal C_{54}
  \longrightarrow H_{54}\longrightarrow1.             \tag{D.42.4}
\]
Every finite level has order $54\cdot3^{10n}$.
The cusp $t=sr$ has order $6\cdot3^n$ there.
The precise finite-base identification, including the
images of $s,r$, is given in the proof.

Consequently the actual simultaneous source image is
\[
 \mathcal Z_{42}
   :=\widehat B/(R_W\cap R_{J,\mathrm{nil}})
     \simeq
     \mathcal W\mathbin{\times_{\mathcal C_{54}}}
                         \mathcal P_{J,\mathrm{nil}}.  \tag{D.42.5}
\]
The map $(\alpha_{41},1_{\mathcal P_{J,\mathrm{nil}}})$
is an actual continuous automorphism of this group.
It preserves the indicated peripheral classes with
multiplier one, fixes $t$, and has a commutator
parameter satisfying both hexagons. Its conjugators
are constructed below from the recorded D.39
conjugators. Thus the specified remaining
$\operatorname{PGL}_2(\mathbb F_7)$ coordinate is now
the identity together with \emph{every} pro-$p$
refinement of its rank-$57$ free kernel, simultaneously.
All D.39--D.41 compact-cover and fixed-fibre CM data
are retained.

The new projection kernels are explicit. Put
$J^+=\operatorname{PSL}_2(\mathbb F_7)$ and
$M=N_J\cap L_{54}$. The joint map to
$J\times_{C_2}H_{54}$ is onto, has order $9072$, and
\[
 [N_J:M]=27,\qquad \operatorname{rank}M=1513.
\]
There is a specified surjection
\[
 \omega:M^{(3)}\longrightarrow\mathbb Z_3^{10},\qquad
 \kappa=\ker\omega,                                    \tag{D.42.6}
\]
whose $10$ by $1513$ integral homology matrix is
supplied by the program. It is the graph projection
induced by $M\subset L_{54}$. With
$K_{\mathrm{new}}=\ker(\mathcal Z_{42}\to\mathcal W)$,
\[
 1\longrightarrow
       \left(\prod_{p\ne3}N_J^{(p)}\right)\times\kappa
  \longrightarrow K_{\mathrm{new}}
  \longrightarrow J^+\longrightarrow1.               \tag{D.42.7}
\]
Moreover, with $K_{T,C}=\ker(\mathcal T\to\mathcal C_{54})$,
\[
 \ker(\mathcal Z_{42}\to\mathcal P_{J,\mathrm{nil}})
    =K_{T,C}\times\mathcal R'\times
                  \prod_{M'\in\mathscr S}A_{M'}.       \tag{D.42.8}
\]
Here $\mathcal R'$ is the closed commutator subgroup
specified in D.40. The group $K_{T,C}$ has the
following two exact descriptions of its successive
kernels:
\[
 \begin{split}
 1&\longrightarrow \mathcal K_2
   \longrightarrow K_{T,C}
   \longrightarrow K_{H,C}\longrightarrow1,\\
 1&\longrightarrow
    \left(\prod_{p\ne3}\mathbb Z_p^{1153}\right)
          \times\ker(P_3)
   \longrightarrow K_{H,C}
   \longrightarrow O_2(H)\longrightarrow1,\\
 P_3&:\mathbb Z_3^{1153}\twoheadrightarrow
                      \mathbb Z_3^{10},\qquad
 \operatorname{rank}_{\mathbb Z_3}\ker(P_3)=1143.
 \end{split}                                          \tag{D.42.9}
\]
Here $\mathcal K_2=\ker(T_3\to C_8)$ is the exact
finite group from D.39, of order $2^{11151}$,
with its indicated source maps. The map $P_3$ is the
stored $10$ by $1153$ integral graph-projection
matrix, completed at $3$; $K_{H,C}$ is
$\ker(\mathcal H\to\mathcal C_{54})$.

An explicit nontrivial kernel still remains in the
full source. Define the actual marked map
\[
 \begin{split}
 \sigma_{11}:\widehat B&\twoheadrightarrow
                 J_{11}=\operatorname{PGL}_2(\mathbb F_{11}),\\
 s&\longmapsto
       \left[\begin{pmatrix}1&0\\0&-1\end{pmatrix}\right],
 \qquad
 r\longmapsto
       \left[\begin{pmatrix}1&-1\\1&0\end{pmatrix}\right].
 \end{split}                                          \tag{D.42.10}
\]
Then
\[
 \sigma_{11}(R_W\cap R_{J,\mathrm{nil}})
       =\operatorname{PSL}_2(\mathbb F_{11}).           \tag{D.42.11}
\]
The source $\widehat F$ maps onto the inverse image
of the identity in $S_3$ in (D.42.5), with this same
source kernel. D.45 excludes a full profinite lift of this
recorded action. Formula (D.42.11) records its source
kernel; the prescribed CM pair remains open for a replacement.

This addresses Part D(e), section 51.16, folio
734 verso, scan 174, and D-after-LM2(1), section
51.15, folios 725--727, scans 156--160.
\end{statement}
\begin{proof}
\emph{The relative completions and their solvable quotients.}
The images of $s$ and $r$ in $J$ have exact orders
two and three. The Bass--Serre graph of the kernel
has $336$ edges and $168+112=280$ vertices, so
$N_J$ is free of rank $336-280+1=57$.
Its profinite completion is the open subgroup
$\ker(\widehat B\to J)$. Each $R_{J,p}$ is
characteristic in that subgroup, hence normal in
$\widehat B$. The map from $\widehat N_J$ to every
finite product of the distinct pro-$p$ completions
is onto: a common quotient of a pro-$p$ group and
a product of groups at other primes is trivial.
Finite-product surjectivity and compactness prove
(D.42.2), including the literal fibre-product maps.
In particular every finite quotient retaining
$J$ has a nilpotent normal kernel over $J$.

Write $B_{\mathrm{even}}=\ker(B\to C_2)$ for the
original sign map. Reidemeister--Schreier rewriting
with the two representatives $1,s$ gives
\[
 B_{\mathrm{even}}
   =\langle r,srs\mid r^3=(srs)^3=1\rangle
   =C_3*C_3.                                         \tag{D.42.12}
\]
Let $\mathcal R_3$ be its full relative pro-$3$
completion over the sign quotient:
\[
 \mathcal R_3
  =\widehat B/
      \ker\bigl(\widehat B_{\mathrm{even}}
                   \to B_{\mathrm{even}}^{(3)}\bigr).
                                                               \tag{D.42.13}
\]
The action of $s$ interchanges the two displayed
order-three generators. The inclusion
$N_J\subset B_{\mathrm{even}}$ shows directly
that the map to $\mathcal R_3$ factors through
$\mathcal P_{J,3}$.

Here is the exact quotient argument needed for
the reverse comparison. Let $E$ be a finite
quotient of $\mathcal P_{J,\mathrm{nil}}$
retaining $J$, with nilpotent kernel $V$.
For every quotient $E/U$, its further quotient
by the image of $V$ is
\[
 (E/U)/\operatorname{im}(V)
       =J/\operatorname{im}(U).                       \tag{D.42.14}
\]
D.41 proved that the normal subgroups of $J$
are $1,J^+,J$. If $\operatorname{im}(U)$ is
$J^+$ or $J$, (D.42.14) and the nilpotence of
the image of $V$ make $E/U$ solvable.
Thus every nonsolvable quotient of $E$ has
the actual quotient $J$, obtained by dividing
out the image of $V$. A finite quotient that
does not initially retain $J$ is a quotient
of one that does, by adjoining the $J$
coordinate, so the assertion holds for every
finite quotient of $\mathcal P_{J,\mathrm{nil}}$.

For the solvable images, put $E^+=\operatorname{im}
(B_{\mathrm{even}}\to E)$, the preimage of $J^+$.
Its stable derived subgroup $(E^+)^\infty$
maps onto the perfect group $J^+$. Therefore
\[
 E^+=(E^+)^\infty V.
 \]
Every solvable image of $E^+$ kills
$(E^+)^\infty$ and is consequently an image
of the nilpotent group $V$. It is generated
by the images of $r,srs$, both of order
dividing three. A finite nilpotent group
generated by such elements is a $3$-group,
as its projections to every other Sylow
factor are trivial. Thus every solvable
source image factors (D.42.13).
Conversely $\mathcal R_3$ is prosolvable
and is an actual quotient of
$\mathcal P_{J,3}$. This proves
\[
 (\mathcal P_{J,\mathrm{nil}})^{\mathrm{sol}}
       =\mathcal R_3,\qquad
 (\mathcal P_{J,p})^{\mathrm{sol}}
       =\begin{cases}\mathcal R_3,&p=3,\\C_2,&p\ne3.
         \end{cases}                                 \tag{D.42.15}
\]
The superscript here means the quotient by
the intersection of all open normal subgroups
with solvable quotient; all maps in (D.42.15)
are the specified source maps.

\emph{The actual order-$54$ quotient.}
The program starts with the literal states
$(q,a,b)$ of D.37. Its map is
\[
 (q,a,b)\longmapsto
       \bigl(q\,O_2(Q),\,a\bmod3,\,b\bmod3\bigr).
                                                               \tag{D.42.16}
\]
There are $54$ values and $128$ elements in
each fibre. It checks all $6912$ right-$s$
and right-$r$ equations, so (D.42.16)
is the actual marked quotient map.
The $32$ elements of $O_2(Q)$ are verified
as a normal subgroup, and their inverse
image with $a=b=0\bmod3$ is the indicated
normal $2$-subgroup of $H$. The quotient
has trivial centre; a normal $2$-subgroup
in a group of order $54$ would have order
two and be central. Thus the kernel is
exactly $O_2(H)$.

Here is an explicit identification of the
finite quotient, rather than only its order.
Put $a=r$, $b=srs$, and $c=aba^{-1}b^{-1}$
in $H_{54}$. Its even subgroup is the
order-$27$ group whose coordinates satisfy
\[
 (u,v,w)(u',v',w')
   =(u+u',\,v+v',\,w+w'+uv'),\qquad
 (u,v,w)\longmapsto b^v a^u c^w,                       \tag{D.42.17}
\]
with all coordinates in $\mathbb F_3$.
This gives the exact isomorphism from the
upper-unitriangular three by three matrices
over $\mathbb F_3$. The involution $s$ acts by
\[
 (u,v,w)\longmapsto(v,u,uv-w).                         \tag{D.42.18}
\]
All $27^2$ multiplication equations and all
$27$ involution equations are checked.
The source generator $r$ has coordinate
$(1,0,0)$; $s$ is the displayed involution.
The map to $Q_0$ kills precisely the centre
$\langle c\rangle$. In particular its even
quotient is the actual abelianization
$C_3^2$ of (D.42.12), with the two indicated
generators, so $K=B_{\mathrm{even}}'$ as
marked subgroups.

The Bass--Serre graph for $L_{54}$ has
$54$ edges and $27+18=45$ vertices. Its
ten fundamental integral cycles form the
stored basis, with an identity minor at
the ten non-tree edges. The same projection
on edges sends the $1153$ cycles of $L$
to the specified matrix $P_3$ and sends
the path of $t^6$ to a vector having a
unit coordinate at $3$. This proves all
rank, order and cusp assertions in
(D.42.3)--(D.42.4).

\emph{The maximal common marked quotient.}
D.41 excludes $J$ as a quotient of $\mathcal W$.
By (D.42.14), every common finite quotient
of $\mathcal W$ and $\mathcal P_{J,\mathrm{nil}}$
is therefore solvable. The kernel of
$\mathcal W\to\mathcal T$ is
$\mathcal R'\times\prod A_{M'}$, whose
every finite image is perfect. Hence the
maximal prosolvable quotient of $\mathcal W$
is precisely $\mathcal T$. By (D.42.15)
the required common quotient is the
maximal relative pro-$3$ source quotient
of $\mathcal T$ over the sign.

The projection $\mathcal T\to\mathcal H$
has kernel $\ker(T_3\to C_8)$, a finite
$2$-group of order $2^{11151}$. This
kernel lies in the even subgroup. Its
image in a relative pro-$3$ source
quotient over $C_2$ is trivial.
It remains to compute that quotient
of $\mathcal H=\widehat B/\overline{[L,L]}$.
Here $L\subset L_{54}$ and
$L_{54}/L=O_2(H)$ has order $128$.
In every such target the closed images
of $L_{54}$ and $L$ coincide: their
quotient is both a quotient of this
$2$-group and a pro-$3$ group.
The image of $L$ is abelian, so the
image of $L_{54}$ is abelian. The target
therefore factors through
$\mathcal C_{54}$. Conversely
$\mathcal C_{54}$ is a quotient of
$\mathcal H$ and is relative pro-$3$
over $C_2$, by (D.42.17)--(D.42.18).
This proves maximality with the actual
source maps in both directions.
The closed subdirect-product description,
or equivalently its finite quotient
versions and compactness, now gives
(D.42.5).

\emph{The retained conjugators and the coprime correction.}
Let $A_W,\beta_W$ be D.41's normalized
conjugators, so
\[
 \alpha_{41}(s)=\beta_Ws\beta_W^{-1},\qquad
 \alpha_{41}(r)=A_WrA_W^{-1}.
\]
On the full odd homology component D.39
gave an exact formula over $\mathbb Z[1/2]$,
with raw action $\operatorname{Int}(x^6)$.
Its global gauge has odd component $-6$.
The present program projects those same
stored integral paths and applies the
same gauge. In the ten-cycle basis its
answer is
\[
 A_W\longmapsto1,\qquad
 \beta_W\longmapsto b_C
       =-\tfrac12(e_7+e_8)
          \in\mathbb Z_3^{10}\subset\mathcal C_{54},
 \qquad sb_Cs^{-1}=b_C.                               \tag{D.42.19}
\]
The notation in the middle is additive
in the abelian kernel. In particular
the second image is nonzero. The code
checks the exact full path vector and
its $s$-invariance, not only a finite
reduction. The multiplier of the
normalized action is still one; on
$\mathcal C_{54}$ the action is the
identity. Its $3$-adic gauge is exactly
$-6$, since $x$ lies in the pro-$3$
even subgroup of this quotient.

The actual joint map
\[
 \mathcal P_{J,3}\longrightarrow
        J\times_{C_2}\mathcal C_{54}                  \tag{D.42.20}
\]
is onto. Indeed every solvable quotient
of $J$ factors through its sign $C_2$,
and the displayed common sign is
retained. The kernel of (D.42.20) is
pro-$3$, being contained in
$N_J^{(3)}$. The element $(1,b_C)$
centralizes the displayed involution
$(s_J,s_C)$. It has a lift
$\widetilde\beta_3$ that centralizes
$s$ in $\mathcal P_{J,3}$.

For completeness, this centralizer lifting
uses the following elementary coprime
conjugacy argument. Lift $(1,b_C)$
arbitrarily to $v$. The two involutions
$vsv^{-1}$ and $s$ have the same image
and lie in the preimage of its order-two
subgroup, with pro-$3$ kernel. In a
finite $3$-group kernel, induction on
its centre first conjugates the two
involutions modulo the centre. Their
remaining ratio is a central element
$z$ with $s(z)=z^{-1}$. Conjugation
by $z^{-1/2}$ removes it. The square
root exists uniquely in this finite
abelian odd-order centre. Compactness
over finite quotients gives the same
conjugacy by an element of a pro-$3$
kernel. Multiplying $v$ by that
element gives $\widetilde\beta_3$,
without changing either coordinate
in (D.42.20).

Its $J$ coordinate is one. Thus
\[
 \widetilde\beta_{\mathrm{nil},p}
     =\begin{cases}\widetilde\beta_3,&p=3,\\1,&p\ne3
       \end{cases}                                    \tag{D.42.21}
\]
is a genuine element of the fibre
product (D.42.2), centralizes $s$,
and maps to $b_C$. The pairs
\[
 A_{42}=(A_W,1),\qquad
 \beta_{42}=(\beta_W,\widetilde\beta_{\mathrm{nil}})
                                                               \tag{D.42.22}
\]
therefore belong to the actual group
(D.42.5). They give the stated product
automorphism on both $s$ and $r$.
It is invertible because both factors
are invertible and the common action
is the identity. It fixes $t$.

Both conjugators lie over the identity
in the retained $S_3$ quotient. Their
parameter is
\[
 f_{42}=s\beta_{42}s\beta_{42}^{-1}
        =(f_{39},1_{\mathcal P_{J,\mathrm{nil}}}).
                                                               \tag{D.42.23}
\]
This pair is the image of an actual
element of $\widehat F'$: lift it
first through the proved source
surjection; its $\mathcal W$ component
has zero $\widehat F^{\mathrm{ab}}$
coordinate, and $R_W\subset\widehat F'$
means that $\mathcal W$ retains this
entire abelianization. Thus every
such lift has zero abelianization.
The order-two and order-three source
relations and the fixed cusp give
the two multiplier-one hexagons
exactly as in D.25. This also proves
that the same global gauge has been
retained, rather than replaced at
unrecorded primes.

\emph{The exact kernels and their completed maps.}
The joint finite source image
$J\times_{C_2}H_{54}$ has order
$336\cdot54/2=9072$. The code enumerates
it with the actual two generators.
Its kernel $M$ has $9072$ graph edges,
$4536+3024=7560$ vertices, and rank
$1513$. Projection to the $H_{54}$
graph gives the matrix in (D.42.6).
Its surjectivity holds even integrally:
$L_{54}/M=J^+$ is perfect, hence
$L_{54}=M[L_{54},L_{54}]$.
The completed matrix has rank ten
over $\mathbb Z_3$, also checked
directly modulo three.

The subgroup $M$ has $3$-power index
$27$ in $N_J$. Its full pro-$3$
topology is induced from $N_J$:
the core of any normal finite
$3$-index subgroup of $M$ over the
$27$ cosets is again of $3$-power
index in $N_J$. Consequently
$M^{(3)}$ is exactly the open
subgroup of $N_J^{(3)}$ lying over
the identity in the even order-$27$
group. The kernel of (D.42.20)
is precisely
\[
 \ker\bigl(N_J^{(3)}
                  \to(\mathcal C_{54})_{\mathrm{even}}\bigr)
      =\ker(\omega)=\kappa.                            \tag{D.42.24}
\]
In particular its further abelian
kernel description is
$1\to(M^{(3)})'\to\kappa\to
\mathbb Z_3^{1503}\to1$, where the
last group is the kernel of the
specified completed homology
matrix and the prime denotes
closed commutators.

The kernel of
$\mathcal P_{J,\mathrm{nil}}\to
\mathcal C_{54}$ maps onto $J^+$,
by (D.42.20). Its kernel over $J$
is the product of all the other
prime kernels and (D.42.24).
This proves (D.42.7). The common
quotient retains the full marked
abelianization $C_6$ of $B$:
already $H_{54}\to Q_0$ does so.
Hence setting its coordinate to
one in the fibre-product formula
for $\mathcal W$ yields the direct
product (D.42.8).

For (D.42.9), first project
$\mathcal T$ onto $\mathcal H$;
its exact finite kernel was
specified in D.39. The homology
kernel of $\mathcal H\to H$
is $\widehat{\mathbb Z}^{1153}$.
Its map to the kernel in
(D.42.4) kills every prime other
than three and is exactly $P_3$
at three. The latter is onto.
For example, the graph verifies
rank ten modulo three; alternatively
the index-$128$ quotient makes
coinvariants exact at three.
The kernel of the resulting map
projects onto $O_2(H)$: every
lift of an element of $O_2(H)$
can have its image in
$\mathbb Z_3^{10}$ removed by
this surjection. This proves
both exact sequences and the
rank $1153-10=1143$.

\emph{An evaluated remaining source kernel.}
The second program verifies
(D.42.10) directly. The two
matrices generate all $1320$
projective matrices, their even
subgroup is the square-determinant
group of order $660$, and
$r,srs$ generate that subgroup.
Every conjugacy class in the
even subgroup and the normal
closure of every nonidentity
class representative are
enumerated. Each such closure
has order $660$, proving
simplicity. Its centralizer
in $J_{11}$ is trivial.
Consequently the normal
subgroups of $J_{11}$ are
$1,\operatorname{PSL}_2(\mathbb F_{11}),J_{11}$.
The faithful degree-twelve
permutation action is recorded;
its order is $1320$, not $12!$.
The actual cusp and loop
orders are $10$ and $5$.

The finite-kernel argument of
D.41 excludes $J_{11}$ from
$\mathcal W$ with no new
assumption. In a finite
subdirect image of its
solvable, perfect, and full
symmetric coordinates, the
solvable-coordinate kernel
has image containing the
perfect even subgroup under
any proposed map to $J_{11}$.
Every perfect-coordinate
kernel has image all of
$J_{11}$, because the derived
subgroup surjects onto that
perfect coordinate and the
kernel therefore surjects
onto the sign. Every
symmetric-coordinate kernel
also has image containing
the even subgroup; otherwise
$J_{11}$ would be a quotient
of a full $S_n$, contradicted
by the normal subgroups of
$S_n$ and the order $1320$.
The iterated commutator of
these coordinate kernels
is trivial, whereas its
image contains the nontrivial
perfect even subgroup.
This is a contradiction.

Likewise $J_{11}$ is not a
quotient of
$\mathcal P_{J,\mathrm{nil}}$:
by (D.42.14) it would have
a quotient $J$ of order
$336$, which does not divide
$1320$. A finite subdirect
image of these two groups
cannot map onto $J_{11}$
either. Its two coordinate
kernels commute, but their
images would both contain
the same nonabelian simple
even subgroup, since neither
coordinate admits that
quotient. Thus $\mathcal Z_{42}$
has no $J_{11}$ quotient.

Its source kernel is normal
and is contained in $R_W
\subset\widehat B'$. Its
image under $\sigma_{11}$
is therefore contained in
the simple even subgroup.
That image cannot be
trivial, since triviality
would factor the indicated
epimorphism through
$\mathcal Z_{42}$. It is
the whole even subgroup,
proving (D.42.11).

The two independent runnable
certificates are
\texttt{d\_cm\_relative\_pronilpotent.py}
and \texttt{d\_cm\_relative\_kernel.py}.
The former reads the saved D.37
graph and D.39 exact odd paths,
without importing or rewriting
the preceding executables.
It checks all finite source
maps, all $1153+1513$ integral
cycle projections, and the
nonzero centralizer vector
(D.42.19). The latter uses
only the Python standard
library for its full finite
enumeration. The new source
quotient and all its finite
characteristic refinements
retaining the marked CM data
have the stated actual
actions. No assertion that
these refinements are
cofinal among all quotients
of $\widehat F$ is made.
\end{proof}
